\newif\ifglobal

%%%%%%%%%%%%%%%%%%%%%%%
%%% Compile to view %%%
%%%%%%%%%%%%%%%%%%%%%%%

%%%%%%%%%%%%%%%%%%%%%%%%%%%%%%%%%%%%%%%%%%%%%%%%%%%
%%% Readme:                                     %%%
%%% https://www.overleaf.com/read/bwdjvfjksvjy/ %%%
%%%%%%%%%%%%%%%%%%%%%%%%%%%%%%%%%%%%%%%%%%%%%%%%%%%

% >>> M?: marks a module setting number or important notice

% >>> M4: Set {./relative-path} to external files for this module <<<
\def \modulefiles {./01-BUSARCH}
% To add an external file, use:
% (Replace '---' with actual filename)
% '\input{\modulefiles/tables/---.tex}' for tables
% '\includegraphics{\modulefiles/figures/---}' for figures
% '\subfile{\modulefiles/---}' for register files


% >>> M1: This module is in global project? <<<
% 'YES' - uncomment; 'NO' - comment out
\globaltrue

\ifglobal
    \documentclass[main\_\_CN.tex]{subfiles}
   
\else
    % Font "FandolSong-Regular" used by ctex does not have GSUB table.
    % It causes warnings that do not affect compilation and can be ignored.
    \PassOptionsToPackage{quiet}{fontspec}

    \documentclass[openany, 10pt]{book}
   
    % >>> M2: In ./00-shared/preamble-trm-module, also find
    %         settings for visibility of todo notes, watermark <<< 
    \usepackage{preamble-trm-module}

    % Enable Chinese typesetting
    \usepackage{ctex}

    % >>> M3: Temporary labels in Chinese
    %         you can add your own labels there <<<
    \myexternaldocument{./00-shared/temp-labels-cn}

\fi %\ifglobal


\setcounter{secnumdepth}{4}


% >>> M5: If this module needs any updates in
% './00-shared/preamble-trm-module.sty'
% (include additional package, etc.), contact your module PM
% or create the file './00-shared/changelog.tex' make the
% necessary updates yourself and document them in this file <<<

% Make text left-aligned and allow latex to break pages naturally, comment out for full justification
\raggedright
\raggedbottom


\begin{document}


% Sets language into which pre-defined bits of text are translated
\selectlanguage{chinese}
\setchineseformatting


% Do not delete this block
\ifglobal
% Add the following front matter if
% this module is outside of global project
\else
    \InputIfFileExists{readme}{}{}
    \clearpage
    \listoftodos
    \listofdonetodos
    \tableofcontents
    %\listoftables
    %\listoffigures
\fi


%%%%%%%%%%%%%%%%%%%%%%%%%
%%% TRM Module Begins %%%
%%%%%%%%%%%%%%%%%%%%%%%%%

\hypertarget{busarch}{}
\chapter{总线架构}
\label{mod:busarch}

\chipname{} SoC 采用多层总线矩阵作为通信架构的核心，通过高效互连的结构连接多个主设备和从设备。该架构支持并发访问，有效缓解访问冲突，使高性能和低功耗的 RISC-V 核心、DMA 控制器、调试接口及各类外设能够在竞争最小的情况下同时运行。通过优化数据流与系统性能，即使多个高速外设同时工作，总线矩阵也能够确保系统高效运转。



\chipname{} SoC 的主系统由以下部分组成：

\begin{itemize}
    \item 主设备 (master)：
        \begin{itemize}
            \item 高性能 (HP) 32 位 RISC-V 处理器 0 的 I-bus、D-bus 和 AHB 总线
            \item 高性能 (HP) 32 位 RISC-V 处理器 1 的 I-bus、D-bus 和 AHB 总线
            \item 低功耗 (LP) 32 位 RISC-V 处理器的 I-bus、D-bus 和 AHB 总线
            \item TCM MEM MONITOR
            \item PSRAM MEM MONITOR
            \item TRACE0
            \item TRACE1
            \item GDMA-AHB-HP
            \item GDMA-AHB-LP
            \item GDMA-AXI
            \item USB OTG\_HS
            \item GMAC
            \item SD/MMC
            \item 2D-DMA
        \end{itemize}
    \item 从设备 (slave)：
        \begin{itemize}
            \item HP \& LP 外设
            \item CPU 外设
            \item ROM
            \item LP SRAM
            \item HP SRAM
            \item 外部 RAM
            \item 外部 flash（仅支持读访问）
        \end{itemize}
\end{itemize}

\begin{tiplisting}
HP \& LP 外设和 CPU 外设所包含的具体外设列表，请参阅章节 \ref{mod:sysmem} \textit{\nameref{mod:sysmem}}。
\end{tiplisting}

\chipname{} 的总线架构如下图所示。


\begin{figure}[H]
    \centering
    \includegraphics[width=1.0\textwidth]{\modulefiles/figures/bus_matrix.pdf}
    \begin{minipage}{0.9\textwidth}
    \small
    注：图中交叉点处的圆圈表示对应列的总线可以访问对应行的从设备。\\
    所有主设备对外部 flash 的访问均为只读。\\
    \end{minipage}
    \caption{\chipname{} 总线架构}
    \label{fig:bus-arch}
\end{figure}




图 \ref{fig:bus-arch} \textit{\nameref{fig:bus-arch}} 中各主设备与总线矩阵的连接说明如下：


\begin{itemize}
    \item \textbf{HP CPU I-bus}：连接高性能 (HP) 处理器 0 和处理器 1 的指令总线与总线矩阵，用于处理器取指。可访问对象包括 ROM、HP SRAM 以及外部 RAM 和外部 flash，其中对外部存储器的访问通过 cache 进行。
    \item \textbf{HP CPU D-bus}：连接高性能 (HP) 处理器 0 和处理器 1 的数据总线与总线矩阵，用于处理器的数据加载和存储，以及调试器的存储器访问。可访问对象包括 ROM、HP SRAM 以及外部 RAM 和外部 flash，其中对外部存储器的访问通过 cache 进行。
    \item \textbf{HP CPU AHB}：连接高性能 (HP) 处理器 0 和处理器 1 的 AHB 总线与总线矩阵，用于访问外设或存储器中的数据。可访问对象包括 HP \& LP 外设、CPU 外设、LP SRAM 以及外部 RAM 和外部 flash。其中，对外部 RAM 和外部 flash 的访问用于调试。经该总线访问外部存储器时将绕过 cache，且访问地址会被重映射至通过 cache 访问时所使用的地址段（地址映射的详细配置请参阅章节 \ref{mod:sysmem} \textit{\nameref{mod:sysmem}}）。
    \item \textbf{LP CPU I-bus 和 D-bus}：连接低功耗 (LP) 处理器的指令总线和数据总线与总线矩阵，分别用于取指和数据访问。可访问对象为 LP SRAM。
    \item \textbf{LP CPU AHB}：连接低功耗 (LP) 处理器的 AHB 总线与总线矩阵，用于访问外设或存储器中的数据。可访问对象包括 HP \& LP 外设、CPU 外设、HP SRAM 以及外部 RAM 和外部 flash。
    \item \textbf{TCM MEM MONITOR}：连接 TCM 存储器监测器与总线矩阵，用于将监测数据写入存储器。可访问对象包括 HP SRAM 和外部 RAM。
    \item \textbf{PSRAM MEM MONITOR}：连接 PSRAM 存储器监测器与总线矩阵，用于将监测数据写入存储器。可访问对象包括 HP SRAM 和外部 RAM。
    \item \textbf{TRACE}：连接 TRACE0 和 TRACE1 模块与总线矩阵，用于执行调试操作。两个 TRACE 接口访问通路相同，可访问对象包括 HP SRAM 和外部 RAM。
    \item \textbf{GDMA}：芯片包含 GDMA-AHB-HP、GDMA-AHB-LP 和 GDMA-AXI 三个 DMA 控制器，用于高速数据搬运。
        \begin{itemize}
            \item \textbf{GDMA-AHB-HP}：基于 AHB 总线协议的 DMA 控制器，可访问对象包括 HP SRAM 以及外部 RAM 和外部 flash。
            \item \textbf{GDMA-AHB-LP}：基于 AHB 总线协议的 DMA 控制器，可访问对象包括 LP SRAM、HP SRAM 以及外部 RAM 和外部 flash。
            \item \textbf{GDMA-AXI}：基于 AXI 总线协议的 DMA 控制器，可访问对象包括 HP SRAM 以及外部 RAM 和外部 flash。
        \end{itemize}
        GDMA 通过这些总线进行高速数据传输，包括外设到存储器、存储器到外设以及存储器到存储器的传输。针对 DMA 访问外部存储器的场景，为确保数据吞吐率，建议尽可能设置较长的 burst 长度。
    \item \textbf{USB OTG\_HS}：高速 USB OTG 控制器，其内置 DMA 连接至总线矩阵，用于 USB 2.0 高速数据传输。可访问对象包括 HP SRAM 以及外部 RAM 和外部 flash。
    \item \textbf{GMAC}：以太网 MAC 控制器，其内置 DMA 连接至总线矩阵，用于网络数据传输。可访问对象包括 HP SRAM 以及外部 RAM 和外部 flash。
    \item \textbf{SD/MMC}：SD/MMC 控制器，其内置 DMA 连接至总线矩阵，用于 SD/MMC 卡的数据读写。可访问对象包括 HP SRAM 以及外部 RAM 和外部 flash。
    \item \textbf{2D-DMA}：连接 2D 图形加速器与总线矩阵，用于 2D 图形数据搬运和处理。可访问对象包括 HP SRAM 以及外部 RAM 和外部 flash。
\end{itemize}

\chipname{} 的 cache 仅用于 CPU 对外部存储器的访问。CPU 的 I-bus 和 D-bus 经 cache 访问外部 RAM，而 GDMA、USB OTG\_HS、GMAC、SD/MMC、2D-DMA 等主设备直接访问外部 RAM，不经过 cache。硬件不维护这两条路径之间的数据一致性，因此当 CPU 与这些主设备共享外部 RAM 中的数据时，需由软件保证一致性。

不一致有两个方向：

\begin{itemize}
    \item CPU 写入的数据可能仍保留在 cache 中、尚未写回外部 RAM，此时主设备读到的是旧数据。软件应在主设备读取前，将相应的 cache 内容写回外部 RAM。
    \item 主设备写入外部 RAM 后，cache 中可能仍保留该地址的旧副本，此时 CPU 读到的是旧数据。软件应在 CPU 读取前，作废 cache 中相应的内容。若主设备写入的是待执行的代码，还需作废指令 cache。
\end{itemize}

内部存储器不经 cache 访问，外部 flash 对所有主设备均为只读，两者都不存在上述问题。



\end{document}
